% Chapter 4 Experiments - Table 4
% Generalization performance on ODIR, RFMiD/RIADD, and EDID.
\begin{table}[!htb]
\centering
\caption{不同测试集上的泛化性能比较。比较了基线模型与所提方法在 ODIR 域内泛化（Off-site $\rightarrow$ On-site）以及 ODIR 到 RFMiD、EDID 跨域泛化中的表现，比较的指标均为Final score。其中，Off-site $\rightarrow$ On-site 同时报告了单目与双目结果，跨域泛化部分报告了各数据集对齐三类标签空间后的单目结果；每个方向均在候选检查点中选择测试集表现最优的模型进行报告。总体上，所提方法在各测试方向上均优于基线模型。 Comparison of generalization performance on different test sets. The table compares the baseline and the proposed method on ODIR in-domain generalization (Off-site $\rightarrow$ On-site) and cross-domain generalization from ODIR to RFMiD and EDID, using Final score as the metric. Off-site $\rightarrow$ On-site reports both monocular and binocular results, while cross-domain generalization reports monocular results after aligning each dataset to a three-class label space. For each direction, the checkpoint with the best test-set performance is selected from candidate checkpoints. Overall, the proposed method outperforms the baseline in all evaluation directions.}
\label{tab:generalization}
\resizebox{\textwidth}{!}{
\begin{tabular}{lcccccccc}
\toprule
\multirow{3}{*}{Model} & \multicolumn{4}{c}{Off-site $\rightarrow$ On-site} & \multicolumn{2}{c}{Off-site $\rightarrow$ RFMiD} & \multicolumn{2}{c}{Off-site $\rightarrow$ EDID} \\
\cmidrule(lr){2-5}\cmidrule(lr){6-7}\cmidrule(lr){8-9}
& \multicolumn{2}{c}{Off-site} & \multicolumn{2}{c}{On-site} & Off-site & RFMiD & Off-site & EDID \\
\cmidrule(lr){2-3}\cmidrule(lr){4-5}\cmidrule(lr){6-7}\cmidrule(lr){8-9}
& Monocular & Binocular & Monocular & Binocular & Monocular & Monocular & Monocular & Monocular \\
\midrule
ResNet50 & 79.37 & 76.19 & 77.30 & 75.23 & 94.39 & 83.86 & 94.30 & 52.16 \\
\textbf{Ours} & \textbf{81.71} & \textbf{78.79} & \textbf{78.98} & \textbf{77.49} & \textbf{95.94} & \textbf{85.06} & \textbf{96.03} & \textbf{55.35} \\
\bottomrule
\end{tabular}
}
\end{table}
